\section{Requisitos no funcionales}

La siguiente tabla recoge los requisitos no funcionales de la plataforma Colivi, relacionados con las características de calidad, seguridad y funcionamiento que debe cumplir el sistema.

\begin{table}[H]
    \centering
    \begin{tabular}{|p{2cm}|p{12cm}|}
        \hline
        \textbf{Código} & \textbf{Nombre} \\ \hline
        RNF-01 & Seguridad y control de acceso \\ \hline
        RNF-02 & Privacidad y protección de datos \\ \hline
        RNF-03 & Rendimiento y tiempos de respuesta \\ \hline
        RNF-04 & Disponibilidad del sistema \\ \hline
        RNF-05 & Usabilidad y accesibilidad \\ \hline
        RNF-06 & Compatibilidad y diseño adaptable \\ \hline
        RNF-07 & Fiabilidad e integridad de la información \\ \hline
        RNF-08 & Mantenibilidad y modularidad \\ \hline
    \end{tabular}
    \caption{Relación de requisitos no funcionales de Colivi}
    \label{tab:requisitos-no-funcionales}
\end{table}

\subsection{Descripción de los requisitos no funcionales}

\begin{description}

    \item[RNF-01.] El sistema deberá garantizar la seguridad mediante autenticación, autorización basada en roles y protección de los recursos frente a accesos no autorizados.

    \item[RNF-02.] El sistema deberá garantizar la privacidad de los usuarios, restringiendo el acceso a la información personal y a los datos privados de cada convivencia.

    \item[RNF-03.] Las operaciones habituales del sistema deberán ejecutarse en un tiempo de respuesta inferior a 3 segundos en condiciones normales de carga, excluyendo los tiempos de respuesta de servicios externos.

    \item[RNF-04.] El sistema deberá mantener una disponibilidad mensual objetivo de, al menos, el 99\%, excluyendo los periodos de mantenimiento planificado.

    \item[RNF-05.] La interfaz deberá ser intuitiva y facilitar la navegación, proporcionando información clara sobre las acciones realizadas y los errores producidos.

    \item[RNF-06.] La aplicación deberá ser compatible con los principales navegadores web y adaptarse a diferentes tamaños de pantalla, tanto en dispositivos móviles como en equipos de escritorio.

    \item[RNF-07.] El sistema deberá garantizar la integridad y consistencia de la información, evitando estados incoherentes o registros incompletos ante errores durante las operaciones.

    \item[RNF-08.] La arquitectura y el código deberán mantener una estructura modular y organizada que facilite el mantenimiento, la detección de errores y la incorporación de nuevas funcionalidades.

\end{description}
